\documentclass[12pt]{article}
\usepackage{palestra1}
\newcommand{\poc}[3]{\begin{tabular}[b]{@{}c@{}}\bebe{#1}{#2}\\[1mm]#3\end{tabular}}
\newcommand{\linea}[1]{\hspace*{4mm}$= #1$\hfill\tikz[baseline=-.6ex]\draw[tinta!45,line width=.5pt](0,0)circle(1.8mm);\\[1.4mm]}
\begin{document}
\begin{ficha}{Potencias}{Su propio frasco}{4}

\begin{encargo}
Elevar al cuadrado es multiplicar un número por sí mismo. En la palestra, el
tirador \textbf{se bebe su propio frasco}.
\end{encargo}

\bloque{1}{Mira cómo se hace}{En el campo}\quad{\small así se juega $(-3)^2$}

\vspace{1.5mm}
\begin{minipage}[c]{.7\linewidth}
\paso{1}\textbf{¿Quién es la base?}\ El \nm{-3}, con su paréntesis.\\[1.5mm]
\paso{2}\textbf{Se bebe su frasco:}\ $(-3)\cdot(-3)$. Fuerza 9.\\[1.5mm]
\paso{3}\textbf{Dos menos:}\ vuelve a su bando. Sale el \nm{9}.
\end{minipage}%
\begin{minipage}[c]{.3\linewidth}\centering
\bebe{-3}{\traiciontriple}
\end{minipage}

\vspace{1.5mm}
\begin{listillon}
Si la base es \textbf{negativa}, la potencia es \textbf{positiva} con exponente
\textbf{par} ($2$, $4$…) y \textbf{negativa} con exponente \textbf{impar} ($3$,
$5$…).
\end{listillon}

\bloque{2}{Ahora tú}{En el campo}\quad{\small ¿con qué número sale a tirar?}

\vspace{1.5mm}
\begin{tabular}{@{}c@{\hspace{10mm}}c@{\hspace{10mm}}c@{\hspace{10mm}}c@{}}
\poc{-2}{\traiciondoble}{\ej{a}$(-2)^2 = $\casillita[7mm]} &
\poc{3}{\triplon}{\ej{b}$(+3)^2 = $\casillita[7mm]} &
\poc{-1}{\traicion}{\ej{c}$(-1)^2 = $\casillita[7mm]} &
\poc{-2}{\traiciondoble\traiciondoble}{\ej{d}$(-2)^3 = $\casillita[7mm]}\\
\end{tabular}

\bloque{3}{¿Con o sin paréntesis?}{Sobre el papel}

\vspace{1mm}
{\renewcommand{\arraystretch}{2.1}
\begin{tabular}{@{}l@{\hspace{14mm}}l@{}}
\ej{a}$(-2)^2 = $ \casillita[8mm] & $-2^2 = $ \casillita[8mm]\\
\ej{b}$(-1)^2 = $ \casillita[8mm] & $-1^2 = $ \casillita[8mm]\\
\end{tabular}}

\alto{$-2^2$ no es $+4$: el cuadrado es del 2, y el menos va fuera. Da $-4$.}

\reto{¿Cuánto vale $(-1)^3$? ¿Y $(-1)^4$? ¿Y $(-1)^{10}$?}

\comprueba{Cada código abre esa cuenta en Practicar.}{%
\qr{1}{j=potencias\&m=7\&c=n-3_e2}\hfill\qr{2a}{j=potencias\&m=7\&c=n-2_e2}\hfill
\qr{2b}{j=potencias\&m=7\&c=n3_e2}\hfill\qr{2d}{j=potencias\&m=7\&c=n-2_e3}\hfill
\qr{3a}{j=potencias\&t=1\&c=r_n2_e2}}

\end{ficha}

% ─────────────────────────────── REVERSO ──────────────────────────────────
\begin{dorso}{Potencias}{Todo junto}{4}

\bloque{1}{Mira cómo se hace}{Sobre el papel}\quad{\small así se resuelve $5 - (-2)^2$}

\vspace{1.5mm}
\paso{1}\textbf{Primero la potencia:}\ $(-2)^2 = +4$.\\[1.5mm]
\paso{2}\textbf{Luego, la cuenta:}\ $5 - (+4)$: se retira un $+4$.\\[1.5mm]
\paso{3}\textbf{La bandera:}\ $5 - 4 = +1$.

\bloque{2}{Ahora tú}{Sobre el papel}\quad{\small primero la potencia}

\vspace{1mm}
{\renewcommand{\arraystretch}{2.2}
\begin{tabular}{@{}l@{\hspace{14mm}}l@{}}
\ej{a}$3 + (-2)^2 = $ \casillita[8mm] & \ej{b}$(-3)^2 + (-4) = $ \casillita[8mm]\\
\ej{c}$1 + (-2)^3 = $ \casillita[8mm] & \ej{d}$-2^2 + 6 = $ \casillita[8mm]\\
\end{tabular}}

\bloque{3}{El ojo de Argos}{Sobre el papel}\quad{\small Listillón se ha equivocado
en \textbf{una} línea: márcala}

\vspace{1.5mm}
\begin{minipage}[t]{.45\linewidth}
\ej{a}$(-3)^2 + 1$\\[1.4mm]
\linea{-9 + 1}
\linea{-8}
\end{minipage}\hfill
\begin{minipage}[t]{.45\linewidth}
\ej{b}$2 - (-1)^2$\\[1.4mm]
\linea{2 - (+1)}
\linea{2 + 1}
\linea{3}
\end{minipage}

\vspace{3mm}
\begin{semaforo}
\sfila{Sé cuál es la base de una potencia}
\sfila{Pongo bien el signo de una potencia}
\sfila{Hago primero las potencias y luego sumo}
\end{semaforo}

\comprueba{Cada código abre esa cuenta en Practicar.}{%
\qr{1}{j=potencias\&m=8\&c=n5_r_n-2_e2}\hfill\qr{2a}{j=potencias\&m=8\&c=n3_s_n-2_e2}\hfill
\qr{2b}{j=potencias\&m=8\&c=n-3_e2_s_n-4}\hfill\qr{2c}{j=potencias\&m=8\&c=n1_s_n-2_e3}\hfill
\qr{2d}{j=potencias\&m=8\&c=r_n2_e2_s_n6}}
\end{dorso}
\end{document}
